% sections/abstract.tex -- abstract. Every number is a \Num macro
% (generated/numbers.tex).
\begin{abstract}
We maintain the Betti numbers over $\Z$ of a simplicial complex of dimension at
most three under insertions and deletions of simplices. Version~1 of this paper
proposed a localised update of a boundary matrix reduced by discrete Morse
theory; that update is not exact, and this version replaces it. For each
dimension we maintain a reduced matrix $R=BV$ together with the change of
basis~$V$. An inserted simplex is reduced against the existing pivots; a deleted
simplex is removed from every basis chain that contains it, by adding either a
cycle or the column with the smallest pivot, which is the rank-only form of known
updates for zigzag persistence and warm-started reductions. We prove that the
Betti numbers are exact after every edit, for every order of the cells, and that
the cost of an edit admits no bound in terms of its size. Discrete Morse theory
enters only through the order: in the order given by a collapse that declares
critical cells of the highest dimension, every gradient pair is an apparent pair
and a build performs column additions only in critical columns. On five workload
families with up to \NumSizeMax{} simplices, every Betti vector reported by an
exact method was compared with at least one reference computation and
\NumCoverageExternalPct\,\% with one outside our engine (closed forms or GUDHI);
none disagreed. Where edits
create and destroy only small cycles, in refined and carved 3-balls and in flag
complexes, the latency of the dynamic reduced basis depends little on the size of
the complex: at \NumWTwoNThirtyThreeSEightSize{} simplices all its variants are
\NumWTwoNThirtyThreeSEightDrbSpeedupFullIdMin--\NumWTwoNThirtyThreeSEightDrbSpeedupFullIdMax{}
times faster than static recomputation, with or without the Morse order. The
Morse order saves nothing at build time against the insertion orders of our
generators, although any structured order gives a much sparser basis than a
random one, and helps between rebuilds only on some workloads. On punctured surfaces the method is at most
\NumWOneSphereLFiveMmhmSpeedupFullId{} times faster on average and has the worse
tail latency, and combinatorial counters are far cheaper; on non-locality
witnesses its cost grows linearly with the complex, as the theory predicts. An
ablation that applies version~1's re-reduction rule without basis tracking (and
without its recompression and gates) errs within the first
\NumLcrFirstWrongMaxAll{} updates of every run on surfaces, 3-balls and flag
complexes; with version~1's periodic recompression (the periodic trigger only, without
its locality trigger or gates) it is still wrong after
\NumLcrRcWFourPlaneNFourThousandDTwoErrPct--\NumLcrRcErrMax\,\% of the updates
on surfaces, cavities and flag complexes.
\end{abstract}
